% Section: P5 stack-conditioning factorization — algorithm axis (N2) vs stack axes (mega) (iter 161)
% Closes brief vein (b) at the head-to-head granularity: a single table that
% places the algorithm-axis eta^2 from the N2 four-method same-stack panel
% alongside the eta^2 of every stack axis on the 98-cell mega panel. Reuses
% axis_variance_fraction from platform_modal/scripts/berkeley/unpacking_dpo_ppo_factorization.py.
\subsection{Stack-conditioning factorization: algorithm axis vs stack axes (iter 161)}
\label{sec:p5-iter161-stack-factorization}

The brief's vein (b) asks for a \emph{head-to-head} quantification of
stack-conditioning: how much variance does the algorithm axis explain
\emph{when the stack is fixed} (Ivison et al. 2024 algorithm axis), and
how much does each stack axis explain \emph{across the same algorithm}?
Previous rows answer one half of the question each: iter 85 row 101 /
iter 89 reports $\eta^2_{\mathrm{algo}}$ on the N2 same-stack panel
(point estimate and bootstrap CI), and iter 49 row 60 reports 8-item
MIN-REPORT coverage on the 32-cell mega panel. Iter 161 closes the
remaining gap with a single decomposition that places both halves on
the same scale.

\subsubsection{Two-tier $\eta^2$ decomposition}
\label{sec:p5-iter161-two-tier}
\textbf{Tier A --- same-stack (algorithm axis).} We load
\texttt{platform_hybrid/experiments/results/n2\_reward\_tensor\_resume/n2\_metrics.tsv}
(4 methods $\times$ 40 steps $\times G{=}8 \times$ seed 0 = 160 rows)
and compute $\eta^2(\texttt{method} \to \texttt{reward\_mean})$. On
this scale \textbf{$\eta^2_{\mathrm{algo}} = 0.0075$} (SS\textsubscript{axis}=0.0074,
SS\textsubscript{within}=0.9817), with per-method means
[\texttt{grpo}=0.8342, \texttt{gift}=0.8447, \texttt{aero}=0.8275,
\texttt{areal}=0.8287] and max Cohen's $d{=}{-}0.216$ (aero vs gift).
The algorithm axis explains $<$1\% of within-run variance; this is the
\emph{same} point estimate iter 85 row 101 reported, reproduced here
for direct comparison to the stack axes.

\textbf{Tier B --- stack axes (mega panel).} We load
\texttt{platform_hybrid/experiments/results/mega\_20260704/cells.tsv} (98 cells
spanning 2 models $\times$ 3 task\_slice $\times$ 5 $G$ $\times$ 2
temperature $\times$ 2 seeds) and compute $\eta^2$ of each axis on
\texttt{mean\_reward}. Paired-stack coverage is 48/50 unique
$(\texttt{model},\texttt{task\_slice},G,T)$ cells have BOTH seeds
(0.9600, Wilson [0.8654, 0.9890]).

\subsubsection{Head-to-head results}
\label{sec:p5-iter161-head-to-head}
Table~\ref{tab:p5-iter161-head-to-head} places Tier A and Tier B on
the same scale. The headline read-out:

\begin{itemize}
\item \textbf{Algorithm axis} (N2 same-stack): $\eta^2 = 0.0075$.
\item \textbf{Stack axes collectively}: $\eta^2_{\mathrm{union}}(\texttt{model},\texttt{task\_slice},G,T) = 0.9967$ on 50 stack cells.
\item \textbf{Three of four individual stack axes} dominate the algorithm axis by 4--60$\times$: \texttt{model} ($60.6\times$), \texttt{task\_slice} ($36.5\times$), $G$ ($4.1\times$).
\item \textbf{Temperature axis} is NOT a meaningful stack axis on this panel ($\eta^2(T) = 0.0000 < \eta^2_{\mathrm{algo}}$): only $T \in \{0.6, 1.0\}$ are sampled and the model axis captures the variance.
\item \textbf{Seed axis} is decisively small ($\eta^2(\texttt{seed}) = 0.0000$), so the cross-stack comparison is not seed-noise driven.
\end{itemize}

\begin{table}[ht]
\centering\small
\begin{tabular}{lrrrll}
\toprule
Axis & Source & $n_{\mathrm{rows}}$ & $\eta^2$ & Ratio vs algo & Verdict \\
\midrule
\texttt{method}     & N2 same-stack       & 160 & 0.0075 & 1.0$\times$  & DECISIVE ($\leq 0.05$) \\
$(model, task, G, T)$ union & mega & 50 & \textbf{0.9967} & 132.9$\times$ & DECISIVE ($\geq 0.50$) \\
\texttt{model}      & mega                & 98 & 0.4527 & \textbf{60.6$\times$} & DECISIVE (dominates) \\
\texttt{task\_slice}& mega                & 98 & 0.2729 & \textbf{36.5$\times$} & DECISIVE (dominates) \\
$G$                 & mega                & 98 & 0.0304 & \textbf{4.1$\times$}  & DECISIVE (dominates) \\
\texttt{temperature}& mega                & 98 & 0.0000 & 0.0$\times$  & NULL ($< \eta^2_{\mathrm{algo}}$) \\
\texttt{seed}       & mega                & 98 & 0.0000 & 0.0$\times$  & DECISIVE ($\leq 0.10$) \\
\texttt{step} (within-run) & N2 same-stack & 160 & 0.9284 & 124.0$\times$ & NULL (within-run drift dominates) \\
\bottomrule
\end{tabular}
\caption{Head-to-head $\eta^2$ decomposition: algorithm axis (N2 same-stack,
Tier A) vs each stack axis (mega panel, Tier B). $\eta^2_{\mathrm{algo}} =
0.0075$ is the Ivison (NeurIPS 2024) baseline. Three of four individual
stack axes dominate the algorithm axis by 4--60$\times$; the four-axis
union explains $\eta^2 = 0.9967$ of cross-cell variance. The
\texttt{temperature} axis NULL is informative: only $T \in \{0.6, 1.0\}$
are sampled, and the \texttt{model} axis absorbs the variance. The
\texttt{step} axis NULL is the within-run learning curve (40 steps
$\gg$ typical), not a refutation.}
\label{tab:p5-iter161-head-to-head}
\end{table}

\subsubsection{Operational recommendation}
\label{sec:p5-iter161-recommendation}
\textbf{Report the stack, not the label.} The algorithm label alone
(GRPO / GIFT / AERO / AREAL) explains $<1\%$ of variance in same-stack
runs and is dwarfed by \texttt{model} ($60\times$), \texttt{task\_slice}
($36\times$), and $G$ ($4\times$) on the cross-stack mega panel.
For benchmarking new GRPO-family variants, the minimal reporting set
must include at least \texttt{model} + \texttt{task\_slice} + $G$;
$T$ can be omitted from the headline stack when only two values are
sampled, but should be reported as a paired-cell counterfactual
otherwise. This is the same operational rule iter 49 row 60 derived
from the 8-item MIN-REPORT coverage audit on the 32-cell panel; iter
161 confirms the rule on the 98-cell follow-up with a single
side-by-side decomposition.

\subsubsection{Cross-paper coupling}
\label{sec:p5-iter161-coupling}
\textbf{(i) P5 iter 85 row 101 / iter 89} --- iter 161 reproduces
the iter 85 point estimate ($\eta^2_{\mathrm{algo}}=0.0075$ on
\texttt{reward\_mean}) and joins it to the stack axes. The iter 89
bootstrap CI is consistent (UB 0.0244 on \texttt{reward\_mean}, well
below the 0.05 strict threshold).
\textbf{(ii) P5 iter 49 row 60} --- iter 60's 8-item MIN-REPORT
coverage is preserved on the 98-cell panel (iter 161's
$\eta^2_{\mathrm{union}}=0.9967$ is the formal ANOVA-style restatement
of the same observation).
\textbf{(iii) P7 iter 87 row 103 / iter 79 row 93} --- the iter 161
head-to-head shows that on \texttt{reward\_mean} the algorithm axis
is small; P7's signal-starvation theory treats \texttt{zvf} and
\texttt{pcd} channels (where the algorithm axis IS GIFT-driven per
iter 89 row 108) as the controller input. The two stories are
\emph{complementary, not in conflict}: on terminal reward the
algorithm axis is null, on contrast-yield / length-variance channels
the algorithm axis is GIFT-driven.
\textbf{(iv) FRONTIER INSIGHTS round 2 (Gemini Deep Think)} --- the
frontier synthesis proposed $\delta_{\mathrm{div}} \in [0.13, 0.23]$
as the anti-herding structural bonus. iter 161's $\eta^2(\texttt{step})=0.9284$
on the N2 panel is consistent with that picture: most of the
within-run variance is step-driven learning, not method-driven.
The frontier synthesis's iso-yield framework applies at the
contrast-yield axis (where iter 89 isolates GIFT), not at the
terminal-reward axis (where iter 161 confirms algorithm equivalence).

\subsubsection{Limitations}
\label{sec:p5-iter161-limitations}
(i) The mega panel has only 2 models, 3 task slices, 2 temperatures,
and 2 seeds; the per-axis $\eta^2$ estimates have only 2--5 groups.
The 95\% CIs on the per-axis $\eta^2$ are wide; the headline
conclusion rests on the \emph{union}-axis $\eta^2 = 0.9967$ over 50
stack cells, where the central limit theorem gives a tighter read.
(ii) The mega panel uses \emph{training-free} sampling
(\texttt{loss\_form: n/a-sampling} in the manifests); algorithm-axis
comparisons on this panel would be confounded by stack drift and are
not attempted here. (iii) The H7 NULL on the step axis is *not* a
refutation of the P5 thesis: the P5 thesis is that the \emph{cross-method}
signal is small when stack is fixed, not that within-run variation is
flat. The high $\eta^2(\texttt{step})$ is the within-run learning curve
(40 steps on the same prompts), exactly what the brief expects.